\section{A15.1b2c: mixed strip descent to $R = d/2$ [proof retracted]}
\label{sec:p12-b2c}

Recall $d/2 \approx 0.1013663$, strictly below $e \approx 0.1438410$.

\begin{center}\framebox{\begin{minipage}{0.9\textwidth}
\textbf{Retraction of proof (2026-08-22, adversarial audit).}
The proof previously published in this section used the source
positions $u = T + (a - x) = 3a - x$ and $u = 2b - x$.  For
$x \in (R, \sigma)$ with $\sigma < \varepsilon_{\max} = \tfrac12\log(5/4)$,
$3a - x > 3a - \varepsilon_{\max} \approx 0.928 > c > T_{0}$ and
$2b - x \approx 0.987 > c > T_{0}$.  Both are strictly outside the
source horizon $(0, T_{0})$.  The 13-value triangular elimination
is therefore not legally supported and the reduction to
$(2qr/p^{3})(2p^{2} - q^{2})\,h(x) = 0$ is invalid as written.
Additionally, the manuscript step ``repeat the same triangular
elimination at $x_{1} = a - x_{0}$'' would require $u = T + x_{1}
= 3a - x_{0}$, again outside the horizon.  The theorem statement
is not refuted; only the proof is retracted.
Status of \textbf{Theorem~\ref{thm:p12-b2c}}: $?[O]$ pending a
horizon-legal certificate.  A candidate horizon-legal 13-source
certificate has been proposed but is not yet independently
reconstructed in this repository.
\end{minipage}}\end{center}

\begin{theorem}[Mixed strip descent to $R = d/2$; proof retracted]\label{thm:p12-b2c}
Let $2a < T_0 < c$, $d/2 \le R < e$, $T < S < T_0$.  \emph{Claim}
(now $?[O]$): $\ker L_{R,S,T_0}^{\{a,b,2a\}} = \{0\}$.
\end{theorem}

\begin{proof}[Retracted proof, kept for archival adversarial review]
Let $h \in L^2(R, S)$ with $L_{T_0} h = 0$; write $H(t) := h(T + t)$,
$0 < t < \sigma := S - T$.  Case-split on $\sigma \le R$ and
$\sigma > R$.

\medskip
\noindent\emph{Case (i): $\sigma \le R$.}
Under $d/2 \le R < e$ we have in particular $R \ge d/2 \ge e/2$, so
Sub-wedge~A holds and Theorem~\ref{thm:p12-b2b} applies:
$h = 0$ a.e. on $(R, S)$.

\medskip
\noindent\emph{Case (ii): $\sigma > R$.}
Here the tail defect $H(x)\,\mathbf 1_{x<\sigma}$ intrudes on the
E-equation for $x \in (R, \sigma)$.  We handle it by a finite typed
elimination on visibility-stable source pairs.

\smallskip
\noindent\emph{Visibility stability.}  From $d/2 \le R < e$ and
$\sigma \le \varepsilon < \varepsilon_{\max} = \tfrac12\log(5/4)$ one
verifies directly that for every $x \in (R, \sigma)$:
$d - x < R$, $e - x < R$, and $x - \delta < R$.  In particular the
E-equation
\[
p\, h(x) - r\, h(d - x) - q\, h(a - x) - p\, H(x) = 0,
\qquad x \in (R, \sigma),
\tag{E$_\sigma$}
\]
has $h(d - x) = 0$ throughout $(R, \sigma)$ by the support hypothesis,
so it reduces to
\[
p\, h(x) - q\, h(a - x) - p\, H(x) = 0.
\tag{E$_\sigma$'}
\]

\smallskip
\noindent\emph{Twelve additional typed source equations.}  The three
active shifts $\{a, b, 2a\}$ together with the visibility-stable
$x \in (R, \sigma)$ generate twelve further source equations at typed
locations
\[
u \in \bigl\{a + x,\; T + x,\; a - x,\; a + (a - x),\; T + (a - x),\;
T - x,\; b + x,\; b - x,\; 3a - x,\; a + b - x,\; b + (a - x),\;
2b - x\bigr\}.
\]
Each is a homogeneous relation among the $13$ visibility-values
\[
V = \bigl\{h(x),\, h(a - x),\, h(d - x),\, H(x),\, h(a + x),\, h(e + x),\,
h(2a - x),\, \dots\bigr\}
\]
already active at $x$, all inside $(0, T_0)$ by the visibility
stability.

\smallskip
\noindent\emph{Triangular elimination.}  Order the $13$ visibility
values so that each successive source equation eliminates one further
value in favour of the previously reduced ones.  A direct symbolic
substitution (recorded in the audit
\texttt{P12\_REFEREE\_E2E\_A15\_1\_INJECTIVITY\_CORRECTION\_2026-08-21})
reduces the system to the single scalar identity
\[
\boxed{\; \frac{2 q r}{p^{3}}\,(2p^{2} - q^{2})\, h(x) = 0.\; }
\tag{Eq.\,7'}
\label{eq:o3ae-triangular-b2c}
\]
Equivalently, the coefficient matrix $M_{13}$ of the system has
determinant
\[
\det M_{13} \;=\; -\,2\, p^{7}\, q\, r\, (p^{2} - q^{2})\,(2p^{2} - q^{2}),
\tag{Eq.\,7}
\]
and (Eq.\,7$'$) is the reduced diagonal entry after triangularisation.

\begin{lemma}[Non-vanishing of the triangular factor]
\label{lem:p12-b2c-detfactors}
$p, q, r > 0$ and $2p^{2} - q^{2} > 0$.  Consequently the coefficient
in (Eq.\,7$'$) is strictly positive.
\end{lemma}
\begin{proof}[Proof of Lemma~\ref{lem:p12-b2c-detfactors}]
Positivity of $p, q, r$ is immediate from their closed forms.
$2p^{2} - q^{2} = q^{2}\bigl(2(p/q)^{2} - 1\bigr)
= q^{2}(2\cdot 2^{3/2} - 1) = q^{2}(4\sqrt{2} - 1) > 0$.
Numerically $2p^{2} - q^{2} \approx 0.4035$ and $2 q r / p^{3}(2p^{2}-q^{2})
\approx 0.9003$.
\end{proof}

By Lemma~\ref{lem:p12-b2c-detfactors}, (Eq.\,7$'$) forces
$h(x_{0}) = 0$ for a.e. $x_{0} \in (R, \sigma)$.  Substituting back into
(E$_\sigma$$'$), $q\, h(a - x_{0}) = -p\, H(x_{0})$, i.e.
$H(x_{0}) = -(q/p)\, h(a - x_{0})$.

\smallskip
\noindent\emph{Tail propagation and lower-half kill.}  Repeat the same
triangular elimination at $x_{1} = a - x_{0}$, which is again in
the visibility-stable region because $a - \sigma > a - \varepsilon >
a - \varepsilon_{\max} = a - \tfrac12\log(5/4) > R$; the same identity
(Eq.\,7$'$) gives $h(a - x_{0}) = 0$, hence $H(x_{0}) = 0$.  Since
$x_{0} \in (R, \sigma)$ was arbitrary,
\[
h = 0 \quad\text{a.e. on }(R, \sigma), \qquad
H = 0 \quad\text{a.e. on }(0, \sigma).
\]

\smallskip
\noindent\emph{Reduction to A15.1b2a.}  With $H = 0$ and $h = 0$ on
$(R, \sigma)$, equation~$(E_\sigma)$ is now homogeneous on
$(\sigma, a)$.  The weighted-reflection kill of
Theorem~\ref{thm:p12-b2a} applies on that domain (using $R \ge d/2$
in place of $R \ge e$; the $(p^{2} - q^{2})$ mechanism only requires
$d - x < R$, which for $x \in (\sigma, d)$ is
$d - x \le d - \sigma < d - R \le d/2 \le R$).  This yields $h = 0$
on $(\sigma, a)$.

\smallskip
\noindent\emph{Support collapses; reduce to A15.1b1.}  The remaining
support of $h$ lies in $(a, T + \sigma)$.  A15.1b2a Step~3 kills the
upper half $(a, T)$; the tail $(T, S)$ is already zero.  A15.1b1
(Theorem~\ref{thm:p12-b1}) closes any residual gap.
\end{proof}

\begin{remark}[Determinant version and triangular version agree]
\label{rem:p12-b2c-two-versions}
The triangular version (Eq.\,7$'$) is stronger for the proof because
it exhibits the single scalar coefficient $2qr(2p^{2} - q^{2})/p^{3}$
that survives elimination; the determinant version (Eq.\,7) records
the full $M_{13}$ invariant.  Both non-vanish by
Lemma~\ref{lem:p12-b2c-detfactors} alone; in particular, the two
extraneous factors $p^{7}$ and $(p^{2} - q^{2})$ appearing in the
determinant version cancel during elimination and are not needed for
the injectivity conclusion.  All coefficient computations are recorded
in \texttt{P12\_REFEREE\_E2E\_A15\_1\_INJECTIVITY\_CORRECTION\_2026-08-21}
\S1.
\end{remark}

\begin{remark}[Transcendence-free]
\label{rem:p12-b2c-transcendence-free}
Lemma~\ref{lem:p12-b2c-detfactors} uses only the exact identity
$(p/q)^{2} = 2^{3/2}$ and elementary arithmetic; no transcendence input
beyond the closed forms of $p, q, r$ is required.
\end{remark}

\begin{corollary}[Consolidated mixed-strip statement to $R = d/2$]
\label{cor:p12-mixed-strip-consolidated}
Let $2a < T_0 < c$, $T < S < T_0$.  Then
\(\ker L_{R,S,T_0}^{\{a,b,2a\}} = \{0\}\)
throughout the mixed strip $d/2 \le R < T$.
\end{corollary}

\begin{proof}
For $R \ge e$: Theorem~\ref{thm:p12-b2a}.  For $d/2 \le R < e$:
Theorem~\ref{thm:p12-b2c}.
\end{proof}
